\section*{Chapitre \ref{ch:mx:empirical-facts-of-order-books} --- Faits empiriques des carnets d'ordres}

\begin{solution}{exo:mx:empirical-facts-of-order-books:1}
Trois ticks (le meilleur cours acheteur moins le prix, 100.00 moins 99.97)~; zéro pour un achat à 100.00, qui rejoint le meilleur.
\end{solution}

\begin{solution}{exo:mx:empirical-facts-of-order-books:2}
$0.96+1.64\times0.9=2.44$ ticks.
\end{solution}

\begin{solution}{exo:mx:empirical-facts-of-order-books:3}
Les \omterm{def:mx:order-types-and-their-uses:hidden}{ordres cachés} sont invisibles tant qu'ils ne sont pas exécutés~: un flux de données ne montre la liquidité cachée qu'à travers ses exécutions, de sorte que la part observable est celle des transactions.
Une part des ordres mesurerait l'intérêt caché, ce que seule la plateforme (ou les données au niveau de l'ordre de son régulateur) peut voir.
\end{solution}

\begin{solution}{exo:mx:empirical-facts-of-order-books:4}
Environ 540. La table de survie estime la distribution du délai avant annulation comme si les exécutions n'avaient pas eu lieu (elles censurent l'horloge)~: c'est la
loi de la patience d'un ordre, et non la part de l'ensemble des ordres annulés aussi vite~; les ordres exécutés tôt auraient été annulés plus tard
dans une proportion égale au risque instantané estimé.
\end{solution}

\begin{solution}{exo:mx:empirical-facts-of-order-books:5}
La hausse des cours des actions a rendu les lots standard de 100 actions coûteux, si bien que davantage d'ordres sont plus petits qu'un lot~; les algorithmes découpent finement les ordres, et les courtiers de détail
acheminent les ordres fractionnés et de petite taille. Un gestionnaire de flux de données ou une vue consolidée qui ignore les lots irréguliers manque la plupart des transactions et, lorsque des lots irréguliers reposent à
de meilleurs prix, le véritable meilleur prix.
\end{solution}

\begin{solution}{exo:mx:empirical-facts-of-order-books:6}
Lorsque la fourchette vaut déjà un tick et qu'elle ne peut, la plupart du temps, se resserrer, les fournisseurs de liquidité se font concurrence par la longueur de la file, et la sélection adverse se manifeste
dans la profondeur et la dynamique des files, au lieu de la fourchette~; la relation vaut pour la fourchette implicite du chapitre~6, et non pour la fourchette affichée.
\end{solution}

\begin{solution}{exo:mx:empirical-facts-of-order-books:7}
La journée simulée~: 1.8~\%, 4.6~\% et 12.4~\% exécutés en 1, 10 et 60 secondes, contre 3.4~\%, 8.7~\% et 15.0~\% pour les grandes actions américaines. Les exécutions dans le
simulateur sont environ deux fois moins rapides à une seconde et proches à une minute, alors que ses annulations sont cent fois plus lentes~: le marché simulé est lent à
réviser ses cotations, et non à négocier, si bien qu'il surestime la durée pendant laquelle un ordre au repos survit à l'arrivée de nouvelles informations.
\end{solution}

\begin{solution}{exo:mx:empirical-facts-of-order-books:8}
Les ordres du simulateur vivent cent fois plus longtemps que les ordres réels et ses fournisseurs de liquidité annulent lentement~: un ordre au repos y fait face à bien moins
de concurrence pour annuler les cotations périmées, et les intervenants plus rapides qui capturent les cotations lentes sont absents. Les taux d'exécution et la sélection adverse doivent être mesurés face à un
marché dont la vitesse d'annulation correspond à celle de la plateforme (chapitre~27), ou en conditions réelles à petite taille.
\end{solution}

\begin{solution}{pb:mx:empirical-facts-of-order-books:1}
\textbf{1.} Un tick derrière le meilleur (1\,508 contre 1\,366 actions dans l'heure du chapitre~1)~; les carnets réels sont aussi en bosse (Bouchaud, Mézard et Potters).
\textbf{2.} 0.090 et 0.044 par action et par seconde~: les ordres périmés au meilleur niveau sont retirés lorsque le prix efficient les traverse.
\textbf{3.} 0.5~\%, 18.6~\% et 40.1~\%.
\textbf{4.} Une loi de puissance d'exposant cumulé proche de 1.5 jusqu'à environ 2\,000 ticks (Londres)~; le simulateur place des ordres sur dix niveaux, avec des taux
décroissant géométriquement, aucun au-delà de neuf ticks.
\textbf{5.} 96.6~\% du temps~; moyenne de 1.042 ticks.
\textbf{6.} $s\approx0.96+1.64\,\sigma_1$, $R^2=0.989$.
\textbf{7.} $R^2=0.48$, pente 0.19~: un carnet à grand tick dont la fourchette ne peut bouger.
\textbf{8.} En forme de U~: 92\,200 actions dans le premier treizième, 24\,700 au plus bas, 55\,600 dans le dernier~; moyenne sur huit séances.
\textbf{9.} Durées allant de l'ajout à chaque annulation ou transaction, horloge redémarrée après les événements partiels, estimées par la méthode de la table de survie, les exécutions censurant
les annulations.
\textbf{10.} Grandes actions 54~\% en moins de 0.1~s et 74~\% en moins de 1~s~; petites actions 37~\% et 49~\%.
\textbf{11.} 0.8~\%, 6.8~\% et 40.0~\%.
\textbf{12.} 11.5~\%.
\textbf{13.} Médiane de 0.19 seconde pour revenir à une fourchette d'un tick~; 92.4~\% en moins d'une seconde (750 événements).
\textbf{14.} 15.5 contre 16.6 en 2026~: proches.
\textbf{15.} \emph{Résultat nommé}~: le simulateur reproduit le profil de profondeur en bosse, le volume intrajournalier en forme de U (implanté), la résilience rapide et le
ratio annulations sur transactions~; il échoue sur la loi de puissance des \omterm{def:mx:empirical-facts-of-order-books:relative}{prix limites relatifs}, sur la relation fourchette--volatilité (grand tick), sur la fourchette à l'ouverture, sur la vitesse d'annulation
(cent fois trop lente) et sur la liquidité cachée (aucune).
\textbf{16.} Des annulations lentes et l'absence de concurrents rapides font paraître les ordres au repos plus sûrs et leurs exécutions plus bénignes~: taux d'exécution biaisés à la hausse, la
sélection adverse à la baisse.
\textbf{17.} Le carnet mince au-delà de neuf ticks~: les gros ordres buteraient sur le vide, surestimant l'impact loin du meilleur niveau.
\textbf{18.} Le nombre d'annulations par transaction est passé de 19.8 à 16.6~; le taux de liquidité cachée est passé de 10.5~\% à 32.2~\%~; le taux de lots irréguliers est passé de 19.2~\% à 68.4~\%.
\textbf{19.} La vitesse d'annulation (révisions rapides des cotations après les mouvements du prix efficient), puis les \omterm{def:mx:order-types-and-their-uses:hidden}{ordres cachés} et un placement plus profond.
\textbf{20.} Un test que tout modèle du marché doit réussir avant que ses réponses ne soient jugées fiables.
\end{solution}

\begin{solution}{iq:mx:empirical-facts-of-order-books:1}
Le meilleur niveau est celui où les exécutions et les annulations de cotations périmées interviennent en premier~; au-delà, les ordres sont hors d'atteinte de la plupart des \omterm{def:mx:the-limit-order-book:orders}{ordres au marché} et s'accumulent,
renouvelés par de nouvelles arrivées. \iqlookfor{flux d'ordres par distance~: taux d'exécution, d'annulation et d'arrivée.}
\end{solution}

\begin{solution}{iq:mx:empirical-facts-of-order-books:2}
C'est la sélection adverse qui le fixe~: si la fourniture de liquidité ne rapporte au plus qu'un profit marginal, la fourchette doit couvrir le mouvement de prix que chaque transaction provoque, et ce
mouvement est la \omterm{def:mx:empirical-facts-of-order-books:vol}{volatilité par transaction}. \iqlookfor{argument du seuil de rentabilité de Wyart et al. et rôle de l'impact des transactions dans la volatilité.}
\end{solution}

\begin{solution}{iq:mx:empirical-facts-of-order-books:3}
Analyse de survie~: traiter l'exécution comme une censure pour l'horloge d'annulation (et réciproquement), estimer avec Kaplan--Meier ou une table de survie, et redémarrer la
horloge après les événements partiels si la question porte sur chaque décision. \iqlookfor{reconnaître la censure~; une moyenne naïve sur les ordres annulés est biaisée.}
\end{solution}

\begin{solution}{iq:mx:empirical-facts-of-order-books:4}
Lister les faits dont dépend le P\&L de la stratégie (dynamique des files, vitesse d'annulation, sélection adverse, résilience), mesurer chacun sur le simulateur et sur
des données réelles ou des statistiques publiées, et ne lancer le backtest qu'une fois que les mesures concordent dans les tolérances~; vérifier que la conclusion résiste aux faits que le simulateur
restitue mal. \iqlookfor{un plan de validation lié à la stratégie, non un score de réalisme générique.}
\end{solution}

\begin{solution}{iq:mx:empirical-facts-of-order-books:5}
L'incertitude est maximale après les nouvelles survenues pendant la nuit~: les fournisseurs de liquidité font face à davantage de flux informé et à davantage de risque de prix, et le carnet se remplit encore après la
enchère d'ouverture~; à mesure que l'information est absorbée, la concurrence resserre la fourchette. \iqlookfor{sélection adverse et risque d'inventaire variant au fil de la journée.}
\end{solution}

\begin{solution}{iq:mx:empirical-facts-of-order-books:6}
Son émetteur cherche (une liquidité cachée ou une contrepartie) ou suit un prix en mouvement, plutôt que d'offrir de la liquidité. Test~: conditionner les mouvements de prix
ultérieurs et les exécutions sur l'activité d'\omterm{def:mx:empirical-facts-of-order-books:life}{ordres fugaces}, et comparer avec les ordres de longue durée~; vérifier si les \omterm{def:mx:empirical-facts-of-order-books:life}{ordres fugaces} se regroupent avant les transactions du même côté.
\iqlookfor{une hypothèse testable avec un groupe de contrôle, et non un simple récit.}
\end{solution}
